\caption{Method comparison for the synthetic VAR(1) benchmark, using the corrected reduced amplification criterion $R_A$.
A fixed family of 29 stable $N=100$ generating operators is used. Each fitted operator is estimated from either one trajectory ($M=1$; faint open symbols) or 25 trajectories ($M=25$; lines). The reference is the optimization-based M2 plane applied to the known generating operator, and therefore tests recovery relative to that operator-dependent reference rather than an independently planted response plane.
(a) Inferred $\widehat R_A$ against the reference $R_{A,\mathrm{ref}}$; the dashed diagonal is equality.
(b) Signed deviation $\widehat R_A-R_{A,\mathrm{ref}}$; the principal panel has a clipped display range, with outlying values indicated at the boundaries and untruncated values retained in the archived CSV. The inset displays the eigenbasis-SVD comparison.
(c) Largest principal angle between the estimated and reference response planes.
(d) Plane-invariant observable alignment on the same held-out trajectory.
The vertical dotted line marks $R_{A,\mathrm{ref}}=1$. Reduced operators with invalid real principal logarithms or unstable spectra are not assigned artificial rate-based ratios; in the $M=1$ data $R_A$ is undefined for 8 M1 and 11 M2 settings. In the companion version, faint dotted colored curves in panels (a) and (b) additionally display $R_\sigma$ and its signed deviation for the $M=25$ results on the same physical operator family, without redefining the horizontal $R_A$ axis. Valid $R_\sigma$ values for $M=1$ are supplied in the accompanying data table.}
